\chapter{Introducing Yourself}
\label{ch:introductions}

Now a first real exchange --- giving your name, asking for someone else's, and
saying you are glad to meet them:

\begin{center}
\mr{माझं नाव मीरा आहे.} \quad(\emph{m\=ajh\d{m} n\=av Mira \=ahe} --- ``My name is Mira.'')\\
\mr{तुमचं नाव काय आहे?} \quad(\emph{tumch\d{m} n\=av k\=ay \=ahe} --- ``What is your name?'')
\end{center}

\emph{Practice lessons: \texttt{lessons/MR-C02-*.md}.}

\section{\mr{नाव} \emph{(n\=av)} --- ``name,'' a word older than Europe}
\label{lesson:naav}

\begin{sounds}
\mr{ना} (\emph{n\=a}) $+$ \mr{व} (\emph{va}) $\to$ \mr{नाव} \emph{n\=av}.
\end{sounds}

\begin{cousinweb}
\mr{नाव} (\emph{n\=av}) $\leftarrow$ Sanskrit \mr{नामन्} (\emph{n\=aman},
``name''). Watch the small Marathi shift: the Sanskrit \emph{-m-} softened to
\emph{-v-}, so \emph{n\=am} became \textbf{n\=av} (Hindi kept the \emph{m}:
\emph{n\=am}). That \emph{n\=aman} goes back to PIE \emph{*h\textsubscript{3}%
n\'{o}mn\d{}} --- the same root as English \textbf{name} and \textbf{noun}, Latin
\emph{n\=omen}, Greek \emph{onoma}. A cousin from before either language existed.
\end{cousinweb}

\begin{grammarlens}
\textbf{A Marathi habit: soft consonants.} \emph{m}$\to$\emph{v} here is one of
many softenings that distinguish Marathi from Hindi. Spot the shift and a
Sanskrit word reads straight into Marathi.
\end{grammarlens}

\section{\mr{माझं} \emph{(m\=ajh\d{m})} --- ``my,'' and gender again}
\label{lesson:majhe}

\begin{sounds}
\mr{मा} (\emph{m\=a}) $+$ \mr{झं} (\emph{jha} with the \emph{anusv\=ara} \mr{ं})
$\to$ \emph{m\=ajh\d{m}}.
\end{sounds}

\begin{cousinweb}
\mr{माझं} (\emph{m\=ajh\d{m}}, ``my'') sits on the first-person root \emph{ma-}
--- the same root behind English \textbf{me}, \textbf{my}, \textbf{mine}, Latin
\emph{meus}, Hindi \emph{mer\=a}.
\end{cousinweb}

\begin{grammarlens}
\textbf{``My'' agrees in three genders.} \mr{माझा} (\emph{m\=ajh\=a}, m.) /
\mr{माझी} (\emph{m\=ajh\=\i}, f.) / \mr{माझं} (\emph{m\=ajh\d{m}}, n.). The noun
\emph{n\=av} (``name'') is \textbf{neuter}, so ``my name'' is \emph{m\=ajh\d{m}
n\=av}. You met this gender idea already in the goodbye \emph{yeto/yete}; it runs
through the whole language.
\end{grammarlens}

\section{\mr{आहे} \emph{(\=ahe)} --- ``is,'' and it comes last}
\label{lesson:aahe}

\begin{sounds}
\mr{आ} (\emph{\=a}, the standalone long-\emph{\=a} vowel) $+$ \mr{हे} (\emph{he})
$\to$ \mr{आहे} \emph{\=ahe}.
\end{sounds}

\begin{cousinweb}
\mr{आहे} (\emph{\=ahe}, ``is'') descends from Sanskrit \mr{अस्ति} (\emph{ásti}),
on the root \textbf{$\surd$as} (``to be''). That \emph{$\surd$as} / PIE
\emph{*h\textsubscript{1}es-} is the ancestor of Latin \emph{est}, German
\emph{ist}, and English \textbf{is} itself --- the oldest verb in the family,
worn to \emph{\=ahe} in Marathi.
\end{cousinweb}

\begin{grammarlens}
\textbf{The verb goes last.} Marathi is subject--object--verb: ``my name Mira
\emph{is}.'' \emph{\=ahe} sits at the end, and it changes with the subject ---
\emph{m\=\i\ \=ahe} (I am), \emph{t\=u \=ahes} (you are), \emph{tumh\=\i\ \=ah\=at}
(you are, resp.), which you will use in the next chapter.
\end{grammarlens}

\section{\mr{माझं नाव \ldots आहे} --- ``my name is\ldots''}
\label{lesson:majhe-naav-aahe}

\mr{माझं} (my) $+$ \mr{नाव} (name) $+$ name $+$ \mr{आहे} (is) $=$ ``my name
is\ldots'' \emph{M\=ajh\d{m} n\=av Arun \=ahe.} The \emph{m\=ajh\d{m}} is neuter
to agree with \emph{n\=av}; the verb comes last.

\section{\mr{तू} / \mr{तुम्ही} \emph{(t\=u / tumh\=\i)} --- ``you,'' familiar and respectful}
\label{lesson:tu-tumhi}

\begin{cousinweb}
\mr{तू} (\emph{t\=u}, familiar ``you'') $\leftarrow$ Sanskrit \mr{त्वम्}
(\emph{tvam}), PIE \emph{*t\=u} --- the root of English archaic \textbf{thou},
Latin \emph{t\=u}, French \emph{tu}, Hindi \emph{t\=u}. \mr{तुम्ही}
(\emph{tumh\=\i}, respectful) is grammatically \textbf{plural} --- honouring one
person as many, exactly like French \emph{vous}; it takes the plural copula
\emph{\=ah\=at}.
\end{cousinweb}

\begin{grammarlens}
\textbf{A choice of respect, every time.} \emph{t\=u} for a friend or child;
\emph{tumh\=\i} for an elder, a stranger, anyone you show courtesy to. Marathi
makes you pick with each address.
\end{grammarlens}

\section{\mr{काय} \emph{(k\=ay)} --- ``what''}
\label{lesson:kaay}

\begin{cousinweb}
\mr{काय} (\emph{k\=ay}, ``what'') grows from the interrogative stem \emph{ka-},
PIE \emph{*k\ensuremath{^w}o-} --- the root behind English \textbf{wh-} words and
Latin \emph{qu-}. Marathi keeps the \emph{k-} plainly across its questions:
\emph{k\=ay} (what), \emph{ko\d{n}} (who), \emph{ku\d{t}he} (where), \emph{k\=a}
(why), \emph{kas\d{m}} (how). Hear that opening \emph{k-} and a question is
coming.
\end{cousinweb}

\section{\mr{तुमचं नाव काय आहे?} --- ``what's your name?''}
\label{lesson:tumche-naav-kaay-aahe}

\mr{तुमचं} (\emph{tumch\d{m}}, ``your,'' neuter to match \emph{n\=av}) $+$
\mr{नाव} $+$ \mr{काय} (what) $+$ \mr{आहे} (is) $=$ ``your name what is?'', verb
last. The ``your'' set mirrors the ``my'' set: \emph{m\=ajh\d{m}} (my) /
\emph{tumch\d{m}} (your, resp.) / \emph{tujh\d{m}} (your, familiar) --- all
neuter here before \emph{n\=av}. To a friend: \mr{तुझं नाव काय आहे?}
(\emph{tujh\d{m} n\=av k\=ay \=ahe?}).

\section{\mr{भेटून आनंद झाला} \emph{(bhe\d{t}\=un \=anand jh\=al\=a)} --- ``a joy to meet you''}
\label{lesson:anand}

\begin{cousinweb}
The warm close is \mr{भेटून आनंद झाला} --- literally ``having met [you], joy
happened.'' \mr{आनंद} (\emph{\=anand}, ``joy, bliss'') is \textbf{Sanskrit} ---
\emph{\=a-} (``fully'') $+$ \emph{nand} (``to rejoice'') --- the \emph{\=ananda}
of \emph{sat-cit-\=ananda} and a common given name. And \emph{bhe\d{t}\=un} is
built on \mr{भेटणे} (\emph{bhe\d{t}\d{n}e}, ``to meet''), the verb behind
Chapter~4's \emph{punh\=a bhe\d{t}\=u}.
\end{cousinweb}

\section{The whole exchange}
\label{lesson:MR-C02-practice}

\begin{center}
\begin{tabular}{@{}ll@{}}
\toprule
Marathi & English \\
\midrule
\mr{माझं नाव मीरा आहे.} & My name is Mira. \\
\mr{तुमचं नाव काय आहे?} & What's your name? \\
\mr{माझं नाव अरुण आहे.} & My name is Arun. \\
\mr{भेटून आनंद झाला.} & Pleased to meet you. \\
\bottomrule
\end{tabular}
\end{center}

Every atom traced: \emph{n\=av} ($\to$ \emph{name}), \emph{m\=ajh\d{m}} ($\to$
\emph{my}), \emph{k\=ay} ($\to$ \emph{what}), \emph{t\=u} ($\to$ \emph{thou}).
And two Marathi distinctives came back: three genders (on \emph{m\=ajh\d{m}}) and
the verb \emph{\=ahe} going \textbf{last}. Next chapter: \emph{tumh\=\i\ kase
\=ah\=at?} (``how are you?'').
